\chapter{Uniform random permutations}
\label{chap:permutations}

\begin{definition}[Uniform random permutation; \han{\S1.1}{3}]
  \label{def:uniform-permutation}
  \fid{adapted}
  For $N\ge0$ let $\SN$ be the set of permutations of $[N]=\{1,\dots,N\}$ (for $N=0$ the set $[0]$ is empty and $\SN$ has one element), with the $\sigma$-algebra of all subsets. The uniform probability measure on $\SN$ gives each permutation mass $1/N!$. A \emph{uniform random permutation} $\pi$ is a random element of $\SN$ with this law; by default, the identity map of $\SN$ under the uniform measure.
\end{definition}

\begin{lemma}[Invariance]
  \label{lem:uniform-invariance}
  \uses{def:uniform-permutation}
  \fid{new}
  If $\pi$ is a uniform random permutation and $\sigma\in\SN$ is fixed, then $\pi^{-1}$, $\pi\circ\sigma$ and $\sigma\circ\pi$ are uniform random permutations.
\end{lemma}
\begin{proof}
  Inversion and composition with a fixed permutation are bijections of $\SN$.
\end{proof}

\begin{proposition}[Marginals; \han{Proposition 1.1(i)}{3}]
  \label{prop:han-1-1-i}
  \uses{def:uniform-permutation}
  \fid{adapted}
  Let $\pi$ be a uniform random permutation, $1\le k\le N$, let $p_1,\dots,p_k\in[N]$ be pairwise distinct positions and $i_1,\dots,i_k\in[N]$. Then
  \[
    \Pr\bigl(\pi(p_1)=i_1,\dots,\pi(p_k)=i_k\bigr)=
    \begin{cases}
      (N-k)!/N! & \text{if } i_1,\dots,i_k \text{ are pairwise distinct},\\
      0 & \text{otherwise.}
    \end{cases}
  \]
\end{proposition}
\begin{proof}
  \uses{lem:uniform-invariance}
  Count the permutations with the prescribed values (Exercise 1.1 of the book). The book states the case $p_j=j$; the general case follows from it by \cref{lem:uniform-invariance}, or directly by the same count.
\end{proof}

\begin{proposition}[\han{Proposition 1.1(ii)}{3}]
  \label{prop:han-1-1-ii}
  \uses{def:uniform-permutation}
  \fid{as printed}
  For $1\le k\le N$ and $q\in[N]$,
  \[
    \Pr\bigl(\pi(1)\le q,\dots,\pi(k)\le q\bigr)=\frac{q(q-1)\cdots(q-k+1)}{N(N-1)\cdots(N-k+1)}\,\ind(q\ge k).
  \]
\end{proposition}
\begin{proof}
  \uses{prop:han-1-1-i}
  Sum \cref{prop:han-1-1-i} over the pairwise distinct values in $[q]$.
\end{proof}

\begin{proposition}[Product moments; \han{Proposition 1.1(iii)--(iv)}{3}]
  \label{prop:han-1-1-products}
  \uses{def:uniform-permutation}
  \fid{as printed}
  Let $\pi$ be a uniform random permutation, $k\in[N]$ and $a\in\R^N$. Then
  \[
    \E\bigl[a_{\pi(1)}a_{\pi(2)}\cdots a_{\pi(k)}\bigr]=\frac{(N-k)!}{N!}\sum_{i_1,\dots,i_k}a_{i_1}a_{i_2}\cdots a_{i_k},
  \]
  where the sum is over the $k$-tuples of pairwise distinct $i_1,\dots,i_k\in[N]$. The case $a_i=i$ is the book's item (iii): $\E[\pi(1)\cdots\pi(k)]=\frac{(N-k)!}{N!}\sum_{i_1,\dots,i_k}i_1\cdots i_k$.
\end{proposition}
\begin{proof}
  \uses{prop:han-1-1-i}
  Sum \cref{prop:han-1-1-i} over the values $i_1,\dots,i_k$.
\end{proof}

\begin{proposition}[Moments; \han{Proposition 1.1(iii)--(iv)}{3}]
  \label{prop:han-1-1-moments}
  \uses{def:uniform-permutation}
  \fid{adapted}
  Let $N\ge2$, $a\in\R^N$, $\bar a=N^{-1}\sum_ka_k$ and $\sigma_a^2=N^{-1}\sum_k(a_k-\bar a)^2$. For positions $i\ne j$,
  \[
    \E a_{\pi(i)}=\bar a,\qquad \Var a_{\pi(i)}=\sigma_a^2,\qquad \Cov\bigl(a_{\pi(i)},a_{\pi(j)}\bigr)=-\frac{\sigma_a^2}{N-1}.
  \]
  In particular ($a_k=k$): $\E\pi(1)=(N+1)/2$, $\Var\pi(1)=(N^2-1)/12$ and $\Cov(\pi(1),\pi(2))=-(N+1)/12$.
\end{proposition}
\begin{proof}
  \uses{prop:han-1-1-i}
  Use the one- and two-dimensional marginals from \cref{prop:han-1-1-i} and $\sum_k(a_k-\bar a)=0$.
\end{proof}
\begin{remark}
  These are the first and second moments that the book lists after its items (iii) and (iv); the general product moments are \cref{prop:han-1-1-products}. The book gives the correlation $-1/(N-1)$ of $a_{\pi(1)}$ and $a_{\pi(2)}$, under the condition $\Var a_{\pi(1)}>0$; the covariance is stated here instead, because it needs no condition, and the correlation follows by dividing by $\sigma_a^2$ when $\sigma_a^2>0$.
\end{remark}

\begin{proposition}[Sampling without replacement; \han{Proposition 1.1(v)--(vi)}{3}]
  \label{prop:han-1-1-sampling}
  \uses{def:uniform-permutation}
  \fid{adapted}
  Let $N\ge2$, $1\le n\le N$, $a,b\in\R^N$, and write $\bar a_{n,\pi}=n^{-1}\sum_{i:\pi(i)\le n}a_i$ for the mean of the sample $\{i:\pi(i)\le n\}$, $S_{ab}=(N-1)^{-1}\sum_k(a_k-\bar a)(b_k-\bar b)$ and $S_a^2=S_{aa}$. Then
  \[
    \E\bar a_{n,\pi}=\bar a,\qquad
    \Cov\bigl(\bar a_{n,\pi},\bar b_{n,\pi}\bigr)=\frac{N-n}{Nn}\,S_{ab},
  \]
  and, if $n\ge2$, the sample covariance $(n-1)^{-1}\sum_{i:\pi(i)\le n}(a_i-\bar a_{n,\pi})(b_i-\bar b_{n,\pi})$ has expectation $S_{ab}$.
\end{proposition}
\begin{proof}
  \uses{prop:han-1-1-moments}
  Expand and use \cref{prop:han-1-1-moments}.
\end{proof}

\begin{proposition}[Maximal entropy; \han{Proposition 1.2}{5}]
  \label{prop:han-1-2}
  \uses{def:uniform-permutation}
  \fid{adapted}
  For a random permutation $\sigma$ with law $p$ on $\SN$, let $H(\sigma)=-\sum_{s\in\SN}p(s)\log p(s)$, with $0\log0=0$. Then $H(\sigma)\le\log N!$, with equality for the uniform law.
\end{proposition}
\begin{proof}
  Gibbs' inequality. The book only states the fact.
\end{proof}

\begin{proposition}[Permutation and independence sampling; \han{Proposition 1.3}{5}]
  \label{prop:han-1-3}
  \uses{def:uniform-permutation}
  \fid{adapted}
  Let $X_1,\dots,X_N$ be i.i.d.\ with law $\mathrm P$ and let $\pi$ be a uniform random permutation independent of $(X_1,\dots,X_N)$. Then for every $n\in[N]$ the variables $X_{\pi^{-1}(1)},\dots,X_{\pi^{-1}(n)}$ are i.i.d.\ with law $\mathrm P$.
\end{proposition}
\begin{proof}
  \uses{lem:uniform-invariance}
  Condition on $\pi$ and use that $(X_{s(1)},\dots,X_{s(N)})$ has the law of $(X_1,\dots,X_N)$ for every fixed $s\in\SN$.
\end{proof}
\begin{remark}
  The independence of $\pi$ and $(X_i)$, which the book's proof uses, is part of the statement here.
\end{remark}

\begin{proposition}[Ranks; \han{Proposition 1.4}{6}]
  \label{prop:han-1-4}
  \uses{def:uniform-permutation}
  \fid{adapted}
  Let $X_1,X_2,\dots$ be i.i.d.\ real random variables with a continuous distribution function $F$, and let $R_i=\sum_{j\le N}\ind(X_j\le X_i)$ for $i\le N$. Then:
  \begin{enumerate}
    \item[(i)] $i\mapsto R_i$ is almost surely a permutation of $[N]$, and it is uniform;
    \item[(ii)] $R_i=N\widehat F_N(X_i)$, where $\widehat F_N$ is the empirical distribution function of $X_1,\dots,X_N$;
    \item[(iii)] for each fixed $i$, $R_i/N-F(X_i)\to0$ almost surely as $N\to\infty$;
    \item[(iv)] $\Cov(R_1/N,F(X_1))=\frac{N-1}{12N}$ and $\Cov(R_1/N,F(X_2))=-\frac1{12N}$ for $N\ge2$.
  \end{enumerate}
\end{proposition}
\begin{proof}
  Exercise 1.2 of the book.
\end{proof}

\section{Examples}

\begin{remark}[Examples 1.1--1.5; \han{\S1.3}{7--8}]
  The book's Examples 1.1--1.5 describe statistics and settings built from a uniform random permutation. They are background for the later chapters and are not formalized.
  \begin{itemize}
    \item Example 1.1 (p.~7): measures of disarray between two permutations $\pi$ and $\sigma$ of $[N]$: Spearman's footrule $\sum_i|\pi(i)-\sigma(i)|$, Spearman's rho $\sum_i(\pi(i)-\sigma(i))^2$, Kendall's tau $\sum_{i,j}\operatorname{sign}(\pi(i)-\pi(j))\operatorname{sign}(\sigma(i)-\sigma(j))$, and Chatterjee's rank correlation, which sums $|\pi(j)-\pi(j')|$ over the indices $j,j'$ that are adjacent in the order given by $\sigma$.
    \item Example 1.2 (p.~7): two-sample rank statistics, the Wilcoxon rank sum $\sum_{i\le m}\pi(i)$ and the Mann--Whitney statistic $\sum_{i\le m}\sum_{j\le n}\ind\bigl(\pi(i)<\pi(m+j)\bigr)$.
    \item Example 1.3 (p.~7): permutation tests. Two samples are pooled into $Z_1,\dots,Z_{m+n}$, and the difference of two statistics computed on $\{Z_i:\pi(i)\le m\}$ and $\{Z_i:\pi(i)>m\}$ is compared with the difference computed on the original samples.
    \item Example 1.4 (p.~8): finite-population inference, the sample $\{z_i:\pi(i)\le n\}$ drawn without replacement from a population that may change with $N$; this is the setting of \cref{chap:permutation-processes} (\cref{def:population}).
    \item Example 1.5 (p.~8): design-based causal inference, where a sample drawn without replacement from a finite population with potential outcomes is split at random into treated and control units.
  \end{itemize}
\end{remark}
